\documentclass[12pt]{article}

\usepackage[margin=1in]{geometry}
\usepackage{times}
\usepackage{setspace}
\usepackage{fancyhdr}
\usepackage{ragged2e}
\setstretch{1.5}

\pagestyle{fancy}
\fancyhf{}
\fancyhead[R]{AyurSpace}
\fancyfoot[L]{Pillai HOC College of Engineering \& Technology (Autonomous), B. E. Computer Engineering}
\fancyfoot[R]{\thepage}
\renewcommand{\headrulewidth}{0.4pt}
\renewcommand{\footrulewidth}{0.4pt}

\begin{document}

\begin{center}
  {\Large \textbf{Abstract}}
\end{center}

\vspace{0.5cm}

\justifying

\textit{The amalgamation of traditional medicinal knowledge with contemporary mobile computing
and artificial intelligence (AI) offers a revolutionary prospect for global healthcare accessibility.
Ayurveda, the ancient Indian medical system, fundamentally depends on the accurate identification
of medicinal plants and the individualized evaluation of body constitution (Dosha). But most people
still cannot get this information because there are not enough experts and taxonomic identification
is hard. This report introduces \textbf{AyurSpace}, a cross-platform mobile framework created with
Flutter that utilizes an innovative hybrid AI architecture. The system combines Plant.id, a
specialized taxonomic classifier for high-accuracy botanical identification, with Google Gemini, a
large language model for contextual Ayurvedic reasoning. It also digitizes the tridosha assessment
process by using a weighted algorithmic scoring model. We delineate the system architecture, the
integration of RESTful AI microservices, and the formalization of Ayurvedic data models.
Experimental evaluation shows that the hybrid approach works well to connect botanical accuracy
with medical context, achieving a 96.4\% accuracy in taxonomy and a 99.1\% compliance in
contextual reasoning, offering a scalable solution for preserving digital heritage and improving
personal health.}

\vspace{0.5cm}

\noindent\textbf{Keywords}\textit{---Ayurveda, Mobile Computing, Generative AI, Plant Identification,
Computer Vision, Gemini API, Flutter, Digital Health}

\end{document}
